Una manera d'obtenir la constant elàstica d'una molla és penjar-hi una massa i mesurar-ne el període de les petites oscil·lacions al voltant de la posició d'equilibri. En la gràfica següent hi ha representada la relació entre la massa penjada de la molla i el quadrat del període de les oscil·lacions:

\figura[0.5]{fig1.png}

\begin{apartat}{1}
A partir de la gràfica, calculeu la constant elàstica de la molla. Si l'amplitud de les oscil·lacions fos de $0,10\ \mathrm{m}$, quina seria l'energia cinètica màxima assolida per la massa en l'oscil·lació?
\end{apartat}

\begin{apartat}{1}
Suposem que la constant elàstica de la molla és de $150\ \mathrm{N\,m^{-1}}$, hi pengem una massa d'1,5~kg i la fem oscil·lar amb una amplitud de $0,20\ \mathrm{m}$. Quina és l'acceleració màxima que assoleix? Si submergim tot el conjunt en un recipient ple d'aigua de manera que la massa oscil·la fins a aturar-se a causa del fregament, quin és el treball fet per la força de fregament que ha aturat l'oscil·lació?
\end{apartat}
